% !htbp
\begin{table}[t]
  \centering
  % ============================= Outer container (single prompt)
  \begin{tcolorbox}[
    rounded corners,
    arc=3pt,
    enhanced,
    width=1\textwidth,
    colframe=refpurple!15, % <----- CHANGED
    coltitle=refpurple,
    colback=white,
    title=\textbf{Evaluation Criteria},
    fonttitle=\bfseries\small,
    boxrule=0.6pt,
    left=2pt,
    right=2pt,
    top=0pt,
    bottom=1pt
  ]
% --------------- Part 2: Previous CUDA implement with scores
\begin{tcolorbox}[
      sharp corners,
      colframe=refgreen,
      colback=refgreen!5,
      colbacktitle=refgreen!15,
      coltitle=refgreen,
      boxrule=0.4pt,
      title={\textbf{Default evaluation criteria}},
      fonttitle=\bfseries\small,
      left=0pt,
      right=0pt,
      top=-4pt,
      bottom=-4pt
    ]
      \begin{center}
        \begin{lstlisting}[style=mystyle, frame=none,framerule=0pt, language=Python]
def check_precision(outputs, outputs_new, max_abs_error, max_rel_error):
    # outputs: list of tensors from reference model
    # outputs_new: list of tensors from LLM generated kernel
    outputs = [outputs] if not isinstance(outputs, list) else outputs
    outputs_new = [outputs_new] if not isinstance(outputs_new, list) else outputs_new
    
    all_abs_diff, all_rel_diff = [], []
    is_accurate = True
    
    # Process each output pair
    for out, out_new in zip(outputs, outputs_new):
      abs_diff = torch.abs(out - out_new)
      rel_diff = abs_diff / (torch.abs(out) + 1e-7) # Add epsilon to avoid division by zero
      all_abs_diff.append(abs_diff.view(-1))
      all_rel_diff.append(rel_diff.view(-1))
    
      # Check if any element exceeds both absolute and relative error thresholds
      if ((abs_diff > max_abs_error) & (rel_diff > max_rel_error)).any():
          is_accurate = False
    
    # Combine all differences
    all_abs_diff = torch.cat(all_abs_diff)
    all_rel_diff = torch.cat(all_rel_diff)
    
    return (1 if is_accurate else 0), all_abs_diff, all_rel_diff
\end{lstlisting}
\end{center}
\end{tcolorbox}

\begin{tcolorbox}[
      sharp corners,
      colframe=refgreen,
      colback=refgreen!5,
      colbacktitle=refgreen!15,
      coltitle=refgreen,
      boxrule=0.4pt,
      title={\textbf{Custom evaluation criteria}},
      fonttitle=\bfseries\small,
      left=0pt,
      right=0pt,
      top=-4pt,
      bottom=-5pt
    ]
      \begin{center}
        \begin{lstlisting}[style=mystyle, frame=none,framerule=0pt, language=Python,upquote=true]
def custom_check_precision(param, outputs, outputs_new):
    dtype_str = param.get(`dtype`, `float16`)
    dtype = getattr(torch, dtype_str)
    if dtype == torch.float32:
        return check_precision(outputs, outputs_new, max_abs_error=0.00001, max_rel_error=0.00001)
    else:
        return check_precision(outputs, outputs_new, max_abs_error=0.001, max_rel_error=0.001)
\end{lstlisting}
\end{center}
\end{tcolorbox}
\end{tcolorbox}
\vspace{-10pt}
\caption{Correctness evaluation criteria for LLM-generated kernels, covering both default and custom methods.}
\label{tab:kernel_correctness}
\end{table}



